\documentclass[12pt]{article}
\usepackage{amsmath}
\usepackage{enumitem}
\usepackage{geometry}
\geometry{margin=1in}
\usepackage{hyperref}
\hypersetup{
    colorlinks=true,
    linkcolor=black,
    urlcolor=blue,
    pdfborder={0 0 0}
}
\title{Ujian Tengah Semester – Kriptografi}
\date{}
\begin{document}
\maketitle

\section*{Petunjuk}
\begin{itemize}
    \item Pilih jawaban terbaik untuk setiap soal.
    \item Lingkari huruf pilihan Anda (A, B, C, atau D).
    \item Nilai tiap soal sama besar.
\end{itemize}

\section*{Soal Pilihan Ganda}
\begin{enumerate}[label=\textbf{Soal \arabic*.}, leftmargin=*, itemsep=1.5em]
    \item Prinsip dasar kriptografi mencakup tiga tujuan utama keamanan informasi. Manakah dari pilihan berikut yang \textbf{bukan} merupakan salah satu tujuan tersebut?
    \begin{enumerate}[label=\Alph*.]
        \item Kerahasiaan (Confidentiality)
        \item Integritas (Integrity)
        \item Otentikasi (Authenticity)
        \item Skalabilitas (Scalability)
    \end{enumerate}

    \item Pada cipher klasik Caesar, setiap huruf pada plaintext digeser dengan nilai tetap $k$. Jika $k = 5$ dan huruf ``Z'' dienkripsi, apa huruf cipher yang dihasilkan (gunakan alfabet Inggris 26 huruf)?
    \begin{enumerate}[label=\Alph*.]
        \item E
        \item D
        \item F
        \item A
    \end{enumerate}

    \item Aritmetika modular merupakan fondasi banyak algoritma kriptografi. Nilai $7^{23}\bmod 5$ adalah:
    \begin{enumerate}[label=\Alph*.]
        \item 2
        \item 3
        \item 1
        \item 0
    \end{enumerate}

    \item Manakah pernyataan berikut yang paling tepat mengenai RSA?
    \begin{enumerate}[label=\Alph*.]
        \item Keamanan RSA bergantung pada kesulitan faktorisasi bilangan komposit besar.
        \item RSA menggunakan fungsi hash kriptografis sebagai kunci publik.
        \item RSA tidak memerlukan operasi modular.
        \item RSA hanya dapat mengenkripsi pesan dengan panjang maksimum 128 bit.
    \end{enumerate}

    \item Dalam algoritma Vigenère, kunci yang dipilih harus:
    \begin{enumerate}[label=\Alph*.]
        \item Selalu berupa satu huruf.
        \item Memiliki panjang yang sama dengan plaintext.
        \item Tidak mengulang pola apa pun.
        \item Menggunakan huruf kapital saja.
    \end{enumerate}

    \item Salah satu sifat grup siklik yang penting untuk AES adalah:
    \begin{enumerate}[label=\Alph*.]
        \item Semua elemen memiliki invers.
        \item Operasi grup tidak komutatif.
        \item Hanya ada satu elemen identitas.
        \item Setiap elemen dapat dihasilkan dari pangkat satu generator.
    \end{enumerate}

    \item Pada proses \textit{key expansion} AES-128, berapa banyak kata (word) yang dihasilkan untuk seluruh ronde?
    \begin{enumerate}[label=\Alph*.]
        \item 44
        \item 60
        \item 128
        \item 32
    \end{enumerate}

    \item Probabilitas dan statistik banyak dipakai dalam kriptanalisis. Manakah contoh paling tepat penggunaannya?
    \begin{enumerate}[label=\Alph*.]
        \item Menghitung nilai hash.
        \item Mengestimasi distribusi ciphertext untuk serangan frekuensi.
        \item Membuat kunci publik untuk RSA.
        \item Menggenerasikan nomor acak dalam algoritma DES.
    \end{enumerate}

    \item Pada analisis keamanan cipher klasik, serangan \textit{known‑plaintext} berarti:
    \begin{enumerate}[label=\Alph*.]
        \item Penyerang hanya mengetahui ciphertext.
        \item Penyerang mengetahui sebagian plaintext dan ciphertext yang bersesuaian.
        \item Penyerang dapat memilih plaintext yang akan dienkripsi.
        \item Penyerang memiliki akses ke kunci privat.
    \end{enumerate}

    \item Algoritma \texttt{modular exponentiation} yang efisien biasanya menggunakan teknik:
    \begin{enumerate}[label=\Alph*.]
        \item Karatsuba multiplication
        \item Square‑and‑multiply
        \item Fast Fourier Transform
        \item Euclidean algorithm
    \end{enumerate}

    \item Dalam konteks kriptografi modern, standar AES menggunakan blok berukuran:
    \begin{enumerate}[label=\Alph*.]
        \item 64 bit
        \item 128 bit
        \item 256 bit
        \item 512 bit
    \end{enumerate}

    \item Manakah pernyataan berikut yang paling tepat tentang hubungan antara kriptografi dan CPMK/CPL pada silabus mata kuliah?
    \begin{enumerate}[label=\Alph*.]
        \item Kriptografi hanya terkait dengan CPMK “Pemrograman C”.
        \item Setiap kompetensi (CPL) harus dapat diukur melalui suatu learning outcome kriptografi.
        \item CPMK dan CPL tidak mempengaruhi penyusunan soal ujian.
        \item Hanya topik matematika bilangan yang relevan dengan CPL.
    \end{enumerate}
\end{enumerate}

\newpage
\section*{Kunci Jawaban}
\begin{enumerate}[label=\textbf{Soal \arabic*.}]
    \item D
    \item A
    \item A
    \item A
    \item B
    \item D
    \item A
    \item B
    \item B
    \item B
    \item B
    \item B
\end{enumerate}

\end{document}
